% Lösungen Einheit 1 von 4 – Zählen und Ergebnismengen (zusammenbau v0.1)
\einheitenkopf{Einheit 1 von 4 · Zählen und Ergebnismengen}

\erg{8}{a) 4 Paare: (grau; gelb), (grau; weiß), (grün; gelb), (grün; weiß) \quad b) $3 \cdot 2 = 6$ Kombinationen \quad c) $5 \cdot 7 = 35$ Kombinationen \quad d) 6 Reihenfolgen: Anna–Ben–Cem, Anna–Cem–Ben, Ben–Anna–Cem, Ben–Cem–Anna, Cem–Anna–Ben, Cem–Ben–Anna \quad e) $5 \cdot 4 \cdot 3 \cdot 2 \cdot 1 = 120$ Sitzordnungen \quad f) 953, 359 (größte, kleinste Zahl) \quad g) 4 Ergebnisse: (rot; rot), (rot; blau), (blau; rot), (blau; blau) \quad h) 6 Paare: (1; 4), (2; 4), (3; 4), (4; 4), (5; 4), (6; 4) \quad i) 3 Dreiecke (von 4 Dreierauswahlen liegt ABC auf einer Geraden) \quad j) $3 \cdot 2 \cdot 1 = 6$ Codes (248, 284, 428, 482, 824, 842)}
\erg{9}{a) 3 Ergebnisse: (Kopf; Kopf), (Kopf; Zahl), (Zahl; Kopf)}
\erg{10}{a) Fehler: (Zahl; Kopf) fehlt – (Kopf; Zahl) und (Zahl; Kopf) sind zwei verschiedene Ergebnisse. Richtig sind 4 Ergebnisse. \quad b) Zu jeder Sorte stehen beide Waffeln da, geordnet nach der Sorte – also sind alle Möglichkeiten erfasst. \quad c) 4 Möglichkeiten: (Saft; Apfel), (Saft; Nuss), (Tee; Apfel), (Tee; Nuss) \quad d) $10 \cdot 10 \cdot 10 = 1\,000$ Codes, also höchstens $1\,000$ Sekunden $\approx 17$ Minuten}
